\documentclass[11pt,a4paper]{article}
\usepackage[T1]{fontenc}
\usepackage[utf8]{inputenc}
\usepackage{mathptmx}           % Times New Roman (Nimbus Roman) klonu
\usepackage[a4paper,margin=2.3cm]{geometry}
\usepackage{amsmath}
\usepackage{amssymb}
\usepackage{graphicx}
\usepackage{booktabs}
\usepackage{array}
\usepackage{caption}
\usepackage[hidelinks]{hyperref}
\usepackage{enumitem}
\setlist{nosep,leftmargin=1.4em}
\captionsetup{font=small,labelfont=bf}
\renewcommand{\figurename}{Şekil}
\renewcommand{\tablename}{Tablo}
\setlength{\parskip}{0.45em}
\setlength{\parindent}{0pt}

\newcommand{\reff}[1]{\texttt{#1}}

\begin{document}

\begin{center}
{\Large\bfseries Kapalı Ortamda PPO ile Lidar Tabanlı Drone Keşfi}\\[0.3em]
{\large Pekiştirmeli Öğrenme Dönem Projesi -- Teknik Rapor}\\[0.8em]
{\normalsize Berker Yiğit -- 220202046}\\
{\small Kocaeli Üniversitesi, Bilgisayar Mühendisliği (II. Öğretim)}\\[0.4em]
{\small Algoritma: PPO (Proximal Policy Optimization) \quad\textbar\quad Branch: \texttt{algo/ppo}}\\
{\small Depo: \url{https://github.com/berkeryigit/rl_drone_pathfinding/tree/algo/ppo}}
\end{center}

\hrule
\vspace{0.6em}

\section{Özet}
Bu çalışmada, GPS'in bulunmadığı kapalı bir binada yalnızca 360 derecelik lidar ve
gürültülü odometri kullanan bir drone'un, engellere çarpmadan mümkün olduğunca çok
alanı keşfetmesi problemi pekiştirmeli öğrenme ile çözülmektedir. Politika, PPO
(Proximal Policy Optimization) ile saf Python/numpy üzerinde kurulu hızlı bir iki
boyutlu Gymnasium ortamında eğitilmiştir. Eğitim beş farklı tohum (seed) ile
tekrarlanmış, eğitilmiş politika deterministik koşuda değerlendirilmiş ve sonuçlar
rastgele ile basit sezgisel (heuristic) iki taban politikasına (baseline) karşı
karşılaştırılmıştır. Aynı ortam ve protokol, ekipteki SAC (\texttt{algo/td3})
teslimiyle birebir aynı tutularak algoritmalar arası adil bir karşılaştırma zemini
sağlanmıştır; tek fark algoritmadır. Sonuçlar, PPO ajanının her iki taban politikasını
da belirgin biçimde geçtiğini ve ortalama oda kapsamasını yüksek seviyeye taşıdığını
göstermektedir.

\section{Giriş ve Neden Pekiştirmeli Öğrenme?}
Keşif problemi, ajanın bilmediği bir haritada nereye gideceğine adım adım karar
vermesini gerektirir. Bu problemin pekiştirmeli öğrenme ile çözülmesinin gerekçesi,
rehberin istediği dört ölçütün de sağlanmasıdır:

\textbf{Ardışıklık (eylem geleceği etkiler).} Drone'un bir adımdaki dönüş ve
ilerleme kararı, bir sonraki adımda göreceği lidar manzarasını ve girebileceği
odaları doğrudan belirler. Yanlış bir koridora sapmak, sonraki onlarca adımın
keşif potansiyelini düşürür. Durum geçişi $p(s'\mid s,a)$ eyleme bağlıdır.

\textbf{Gecikmeli ödül.} Yeni bir odaya ulaşmanın getirisi, o odaya götüren kapıyı
bulmak için harcanan boş adımlardan sonra gelir. Altı odanın tamamını gezmenin büyük
bonusu ise ancak bölümün sonuna doğru elde edilir. Ajan, anlık küçük cezaları göze
alıp ileride büyük getiriyi toplamayı öğrenmek zorundadır.

\textbf{Bilinmeyen, büyük durum uzayı.} Gözlem 75 boyutlu sürekli bir vektördür ve
ortam stokastiktir (lidar, odometri ve rüzgar gürültüsü; zamanla yer değiştiren
hareketli engeller). Durum uzayı tablo yöntemleriyle taranamayacak kadar büyüktür;
bu yüzden fonksiyon yaklaşımı (sinir ağı) gerekir.

\textbf{Daha ucuz bir alternatifin bulunmaması.} A* veya BFS gibi klasik arama
yöntemleri, haritanın ve hedefin önceden bilindiği, deterministik ve statik ortamlar
için uygundur. Burada harita önceden bilinmez (ajan onu lidar ile keşfederek çıkarır),
hareketli engeller ortamı zamanla değiştirir ve tek bir hedef yoktur; amaç kapsamayı
en büyüklemektir. Bu koşullarda klasik planlayıcılar doğrudan uygulanamaz; problem
bir sıralı karar (MDP) problemidir ve pekiştirmeli öğrenme doğal çözümdür.

\section{Ortam ve MDP Formülasyonu}
Eğitim ortamı, mevcut Gazebo \reff{multi\_room} dünyasının hızlı bir iki boyutlu
yaklaşımıdır (\reff{env/fast\_2d\_drone\_env.py}). Sekiz metre yarıçaplı, altı odalı
bir bina, $0{,}5$ m çözünürlüklü bir keşif ızgarası ($32\times32$ hücre), 64 ışınlı
bir lidar ve üç adet zamanla yer değiştiren hareketli engel içerir. Aynı dosya,
SAC teslimindeki ortamla bayt-bayt aynıdır (MD5 doğrulanmıştır); böylece iki
algoritma birebir aynı MDP üzerinde karşılaştırılır.

\textbf{Durum (gözlem) $s$.} 75 boyutlu, $[0,1]$ aralığına normalize sürekli bir
vektör:
\begin{itemize}
\item 64 lidar ışını mesafesi (\reff{LIDAR\_MAX}=10 m'ye bölünmüş),
\item ölçülen yönelimin kosinüs ve sinüsü (2),
\item bir önceki eylem $(v_x, v_y, \omega)$ (3),
\item keşif ilerlemesi ve ziyaret edilen oda oranı (2),
\item en yakın engel mesafesi ve boştalık (idle) sayacı (2),
\item en yakın kapıya uzaklık ve duvara yakınlık (2).
\end{itemize}
Yönelim ve poz bilgisine odometri gürültüsü eklenir; yani \emph{ölçülen} durum,
\emph{gerçek} durumdan sapar (Bölüm~\ref{sec:savunma}).

\textbf{Eylem $a$.} Üç boyutlu sürekli bir vektör, \reff{Box}$(-1,1)^3$:
gövde çerçevesinde ileri/geri hız, yanal hız ve dönüş hızı. Komutlar fiziksel
sınırlara ölçeklenir (\reff{V\_MAX}=1.2 m/s, \reff{W\_MAX}=1.5 rad/s) ve üzerine
rüzgar gürültüsü eklenir.

\textbf{Ödül $r$.} Ödül fonksiyonu koddaki \reff{step()} metodunda (\reff{fast\_2d\_drone\_env.py},
158--185.\ satırlar) ile birebir aynıdır:
\begin{align*}
r \;=\; -0{,}01
&\;+\; 3{,}0\cdot \mathbb{1}[\text{yeni voxel}]
\;+\; 20{,}0\cdot \mathbb{1}[\text{yeni oda}]\\
&\;-\; 0{,}10\cdot \mathbb{1}[\text{50 adımdır yeni voxel yok}]\\
&\;-\; 0{,}3\cdot \max\!\Big(0,\tfrac{0{,}7-\ell_{\min}}{0{,}7}\Big)
\;-\; 40{,}0\cdot \mathbb{1}[\text{çarpışma}]\\
&\;+\; 60{,}0\cdot \mathbb{1}[\text{6 odanın tamamı}]
\;-\; 0{,}10\cdot \mathbb{1}[\text{son 60 adımda tekrar ziyaret}]
\end{align*}
Burada $\ell_{\min}$ en küçük lidar mesafesidir. Tasarım felsefesi, toplam getirinin
yaklaşık olarak keşfedilen alanla orantılı olması, çarpışmanın ise ağır
cezalandırılmasıdır. Bölüm, çarpışmada veya altı odanın tamamı gezildiğinde
sonlanır (\emph{terminated}); 600 adıma ulaşınca kesilir (\emph{truncated}).

\section{PPO Algoritması ve Hiperparametreler}
PPO, aktör-kritik ailesinden \textbf{on-policy} bir politika gradyanı yöntemidir.
Aktör, gözlemi alıp eylem üzerinde bir Gauss dağılımı üretir (sürekli eylem); kritik,
durumun değerini tahmin eder. PPO'nun ayırt edici yanı, politika güncellemesini bir
\emph{clipped surrogate} hedefiyle sınırlamasıdır: yeni politikanın eskisine oranı
$1\pm\varepsilon$ (\reff{clip\_range}=0.2) aralığında kırpılır, böylece tek adımda
politikanın aşırı sapması engellenir ve eğitim kararlı kalır. Avantaj tahmini için
GAE (\reff{gae\_lambda}=0.95) kullanılır.

PPO on-policy olduğu için \textbf{replay buffer yoktur}: her rollout'ta toplanan veri,
birkaç epoch boyunca kullanılıp atılır. Bu, SAC/DQN gibi off-policy yöntemlere göre
örnek verimliliğini düşürür; biz bunu sekiz paralel ortam (\reff{num\_envs}=8) ve daha
uzun eğitim ile telafi ettik. Adil karşılaştırma için \reff{learning\_rate}, $\gamma$,
ağ mimarisi ve paralel ortam sayısı SAC teslimiyle aynı tutulmuştur.

\begin{table}[h]
\centering\small
\caption{PPO hiperparametreleri (\reff{config.yaml})}
\begin{tabular}{@{}lll@{}}
\toprule
Parametre & Değer & İşlevi \\
\midrule
\reff{learning\_rate} & $3\times10^{-4}$ & adım büyüklüğü (SAC ile aynı) \\
\reff{n\_steps} & 1024 & ortam başına rollout uzunluğu \\
\reff{batch\_size} & 512 & minibatch boyutu \\
\reff{n\_epochs} & 10 & rollout başına geçiş sayısı \\
\reff{gamma} ($\gamma$) & 0.98 & indirim faktörü (SAC ile aynı) \\
\reff{gae\_lambda} & 0.95 & avantaj yumuşatma \\
\reff{clip\_range} & 0.2 & politika klips oranı \\
\reff{ent\_coef} & 0.005 & entropi (keşif) bonusu \\
\reff{vf\_coef} & 0.5 & kritik kayıp ağırlığı \\
\reff{max\_grad\_norm} & 0.5 & gradyan kırpma \\
\reff{net\_arch} & [256, 256] & ağ genişliği (SAC ile aynı) \\
\reff{num\_envs} & 8 & paralel ortam (SAC ile aynı) \\
\reff{total\_timesteps} & 1.5\,M & seed başına eğitim bütçesi \\
\bottomrule
\end{tabular}
\end{table}

\section{Deney Düzeni}
Eğitim, \reff{seeds.txt}'teki beş tohumla (7, 13, 42, 123, 2025) tekrarlanmıştır.
Her tohum için ayrı bir politika eğitilmiş; eğitim sırasında her 10.000 adımda bir
deterministik eval yapılarak en iyi model saklanmıştır. Eğitim sonunda her politika,
eğitimden ayrı bir RNG akışıyla (tohum$+20000$) 20 bölüm boyunca \textbf{deterministik}
(greedy) koşturularak değerlendirilmiştir; çünkü eğitim eğrisi tek başına kanıt
sayılmaz. Taban politikalar (rastgele ve sezgisel) aynı ortamda 20'şer bölüm
koşturulmuştur. Hiperparametre duyarlılığı için iki parametre üçer değerle taranmıştır
(\reff{learning\_rate}: $10^{-4}, 3\times10^{-4}, 10^{-3}$; \reff{ent\_coef}:
$0{,}0, 0{,}01, 0{,}05$), her biri üç tohumla. Tüm rastgelelik
\reff{numpy.random.default\_rng} ile üretilmiştir; \reff{np.random.seed} ve
\reff{np.random.choice} kullanılmamıştır.

\section{Sonuçlar}
\label{sec:sonuclar}
Beş tohumun deterministik eval özeti Tablo~\ref{tab:ozet}'de verilmiştir. Aşağıdaki
beş grafik, rehberin B bölümü zorunlu standardına uygundur: her grafikte başlık,
birimli eksen etiketi, lejant, beş tohumun ortalaması ile $\pm$ standart sapma bandı
ve hareketli ortalama penceresi bulunur. Öğrenme ve eval eğrilerinde y-ekseni
\textbf{episode getirisidir} (anlık/per-step ödül değildir).

% --- SONUCLAR TABLOSU (egitim sonrasi doldurulacak) ---
\begin{table}[h]
\centering\small
\caption{Deterministik eval özeti (seed $\times$ metrik). Kaynak: \reff{sonuclar/sonuclar.csv}}
\label{tab:ozet}
\begin{tabular}{@{}rrrrrr@{}}
\toprule
seed & n\_ep & eval getiri & eval std & başarı & kapsama \% \\
\midrule
7 & 6105 & 470.08 & 34.28 & 1.00 & 99.75 \\
13 & 6007 & 448.82 & 156.79 & 0.90 & 91.98 \\
42 & 6243 & 455.87 & 24.91 & 1.00 & 99.84 \\
123 & 5917 & 443.37 & 114.05 & 0.95 & 95.91 \\
2025 & 6133 & 474.58 & 113.19 & 0.95 & 95.94 \\
\midrule
\textbf{Ort.} & \textbf{6081} & \textbf{458.5} & -- & \textbf{0.96} & \textbf{96.68} \\
\bottomrule
\end{tabular}

\end{table}

\subsection*{Grafik 1 -- Öğrenme Eğrisi}
\IfFileExists{../sunum/grafikler/1_ogrenme_egrisi.png}{%
\begin{center}\includegraphics[width=0.82\textwidth]{../sunum/grafikler/1_ogrenme_egrisi.png}\end{center}}{}
\textbf{Yorum.} \emph{Gözlem:} Eğitim episode getirisi yaklaşık 290'dan başlayıp 0,5
milyon adımda 470 dolayına tırmanır ve 0,6 milyondan sonra 480 düzeyinde platoya
oturur; tohumlar arası standart sapma bandı zamanla daralır. \emph{Karşılaştırma:}
Beş tohumun tamamı aynı seyri izler ve son üçte birde aralarındaki fark çok küçüktür,
yani sonuç tohuma karşı kararlıdır. \emph{Açıklama:} PPO'nun clipped surrogate hedefi
politikayı her güncellemede sınırlı adımlarla iyileştirdiğinden getiri düzgün ve
sıçramasız artar; plato, ajanın altı odayı düzenli gezmeyi öğrenmesiyle ulaşılan
tavandır. \emph{Sonuç:} 1,5 milyon adım bu ortamda PPO için yakınsamaya fazlasıyla
yeter; ek eğitim getiriyi anlamlı biçimde artırmaz.

\subsection*{Grafik 2 -- Test (Eval) Eğrisi}
\IfFileExists{../sunum/grafikler/2_eval_egrisi.png}{%
\begin{center}\includegraphics[width=0.82\textwidth]{../sunum/grafikler/2_eval_egrisi.png}\end{center}}{}
\textbf{Yorum.} \emph{Gözlem:} Deterministik (greedy) eval getirisi yaklaşık 200'den
0,4 milyon adımda 470 dolayına çıkar ve eğitim eğrisiyle neredeyse aynı platoda kalır.
\emph{Karşılaştırma:} Deterministik eval, stokastik eğitim eğrisiyle aynı seviyeye
oturur; aradaki fark ihmal edilebilir düzeydedir. \emph{Açıklama:} Bu, yüksek
getirinin yalnızca eğitimdeki rastgele keşif sayesinde değil, öğrenilen politikanın
kendisiyle elde edildiğini gösterir; greedy koşuda da başarı korunur.
\emph{Sonuç:} Politika gerçekten genelleşmiştir ve sonuç eğitim eğrisiyle değil,
ayrı bir deterministik koşuyla doğrulanmıştır (rehber B2 şartı).

\subsection*{Grafik 3 -- Loss Eğrisi (Policy Gradient ve Value Loss)}
\IfFileExists{../sunum/grafikler/3_loss_egrisi.png}{%
\begin{center}\includegraphics[width=0.82\textwidth]{../sunum/grafikler/3_loss_egrisi.png}\end{center}}{}
\textbf{Yorum.} \emph{Gözlem:} Value (critic) loss önce 0,3--0,4 milyon adım
civarında 140 dolayında bir tepe yapar, ardından düzenli düşerek 30 seviyesine iner;
policy gradient (actor) loss ise $-0{,}0065$ ile $-0{,}004$ arasında küçük kalıp
yataya oturur. \emph{Karşılaştırma:} Critic loss'un tepe-sonra-düşüş seyri, getirinin
en hızlı yükseldiği döneme denk gelir. \emph{Açıklama:} Ajan büyük ödülleri (yeni
oda, altı-oda bonusu) keşfettikçe değer hedefleri büyür ve critic başta zorlanır
(tepe); politika kararlılaşınca hedefler oturur ve critic loss düşer. Actor loss'un
küçük ve yatay kalması, PPO'nun klips mekanizmasının güncellemeleri sınırlamasından
kaynaklanır. \emph{Sonuç:} Critic'in düşen loss'u öğrenmenin sağlıklı yakınsadığını,
actor loss'un düşük ve yatay olması ise eğitimin kararlı (patlamayan) olduğunu gösterir.

\subsection*{Grafik 4 -- Hiperparametre Duyarlılığı}
\IfFileExists{../sunum/grafikler/4_hiperparametre_duyarlilik.png}{%
\begin{center}\includegraphics[width=0.82\textwidth]{../sunum/grafikler/4_hiperparametre_duyarlilik.png}\end{center}}{}
\textbf{Yorum.} \emph{Gözlem:} \reff{ent\_coef}'te en yüksek ve en az değişken sonuç
0,05'te (368) elde edilir; \reff{learning\_rate}'te en yüksek sonuç $10^{-4}$'te
(375) olup oran büyüdükçe hafifçe düşer ($3\times10^{-4}$: 352, $10^{-3}$: 343). Bu
koşular hiperparametre etkisini izole etmek için kısa bütçeyle (200k adım)
yapıldığından getiriler 1,5 milyonluk ana koşulardan düşüktür. \emph{Karşılaştırma:}
Her iki parametrede de farklar hata bantları içinde kalır; yani ajan, bu aralıkta
hiperparametre seçimine fazla duyarlı değildir. \emph{Açıklama:} Yüksek
\reff{ent\_coef} daha çok keşfe yol açtığından kısa bütçede odaları daha hızlı bulup
biraz daha yüksek getiri verir; çok yüksek \reff{learning\_rate} ($10^{-3}$) ise
güncellemeleri sertleştirip erken kararlılığı hafifçe bozar (keşif--sömürü dengesi).
\emph{Sonuç:} Seçilen taban değerler ($\text{lr}=3\times10^{-4}$,
$\text{ent\_coef}=0{,}005$) güvenli orta noktadır ve sonuç bu aralıkta robust
olduğundan ince ayar kritik değildir.

\subsection*{Grafik 5 -- Baseline Karşılaştırma}
\IfFileExists{../sunum/grafikler/5_baseline_karsilastirma.png}{%
\begin{center}\includegraphics[width=0.7\textwidth]{../sunum/grafikler/5_baseline_karsilastirma.png}\end{center}}{}
\textbf{Yorum.} \emph{Gözlem:} Deterministik eval ortalama getirisi PPO ajanında 459,
sezgisel (heuristic) taban politikada 402, rastgele politikada ise 49'dur.
\emph{Karşılaştırma:} Getiri olarak PPO sezgiseli az farkla geçse de, kapsama
(\%96,7'ye karşı \%79,9) ve çarpışma oranı (0,04'e karşı 0,45) bakımından fark çok
büyüktür; rastgele politika her ölçütte çok geridedir. \emph{Açıklama:} Sezgisel
politika ``ön açıksa ilerle'' kuralıyla bazı odaları gezebilir ama kapıları
planlayamaz ve sık çarpar; PPO ise lidar ve öğrenilen değer fonksiyonu sayesinde hem
daha çok oda gezer hem de neredeyse hiç çarpmaz. \emph{Sonuç:} Ajanın öğrenmesi
gerçek bir kazanımdır; problem basit bir kuralla veya rastgelelikle çözülemez, bu da
pekiştirmeli öğrenme seçimini doğrular (rehber B5).

\section{PPO ve SAC Karşılaştırması}
Ortam, seed listesi ve değerlendirme protokolü ekipteki SAC (\texttt{algo/td3})
teslimiyle birebir aynı tutulduğundan, iki algoritma doğrudan karşılaştırılabilir.
Tablo~\ref{tab:karsilastirma}, her iki algoritmanın beş tohum ortalamasını verir;
SAC değerleri \texttt{algo/td3} branch'inin teslim ettiği \reff{sonuclar.csv}'den
alınmıştır. Tek farklar algoritma ailesi ve eğitim bütçesidir.

\begin{table}[h]
\centering\small
\caption{PPO ve SAC ortalama sonuçları (5 tohum, deterministik eval; aynı ortam ve protokol)}
\label{tab:karsilastirma}
\begin{tabular}{@{}lcc@{}}
\toprule
Ölçüt & PPO (bu çalışma) & SAC (\texttt{algo/td3}) \\
\midrule
Aile & on-policy & off-policy \\
Replay buffer & yok & var \\
Eğitim adımı & 1.5\,M & 500\,k \\
Eğitim episode sayısı & $\sim$6081 & $\sim$1058 \\
\midrule
Ort. eval getirisi & \textbf{458.5} & 321.0 \\
Eval başarı oranı (6 oda) & \textbf{0.96} & 0.27 \\
Ort. kapsama \% & \textbf{96.7} & 92.6 \\
Eval çarpışma oranı & \textbf{0.04} & 0.07 \\
\bottomrule
\end{tabular}
\end{table}

PPO, dört kalite ölçütünün tamamında (getiri, başarı, kapsama, çarpışma) SAC'ın
önündedir; özellikle altı odanın tamamını gezme başarısında fark belirgindir (0,96'ya
karşı 0,27). Ancak bu üstünlük bedelsiz değildir: PPO bu sonuca 1,5 milyon adımda
ulaşırken SAC yalnızca 500 bin adım kullanmıştır. Bunun nedeni doğrudan algoritma
ailesidir; SAC off-policy olduğu için replay buffer'daki her deneyimi defalarca
kullanır ve adım başına daha örnek-verimlidir, PPO ise on-policy olduğundan her
rollout verisini bir kez kullanıp atar. Dolayısıyla tablo, dersin temel ayrımını
somutlaştırır: SAC az adımda makul bir politikaya hızla ulaşır, PPO ise daha çok adım
karşılığında bu ortamda daha yüksek nihai başarıya çıkar. Karşılaştırmanın ortam
tarafı tamamen adildir; eşit adım bütçesinde değerlendirme ileri bir çalışma konusudur.

\section{Savunma Analizi}
\label{sec:savunma}
Bu bölüm, rehberin dört kavram alanını projeye uygulayarak yanıtlar.

\textbf{Deterministik mi, stokastik mi?} Ortam \textbf{stokastiktir}; bu kod
satırlarıyla kanıtlanır. Geçiş gürültüsü: \reff{step()} içinde hız komutlarına
$\mathcal{N}(0,\sigma_{\text{wind}})$ rüzgar eklenir; aynı $(s,a)$ çifti farklı $s'$
verir, yani $p(s'\mid s,a)$ dejenere değildir. Gözlem gürültüsü: lidar
(\reff{\_compute\_lidar}) ve odometri (\reff{\_make\_obs}) ölçümlerine Gauss gürültüsü
katılır. Bu nedenle aynı durumdan aynı eylem dizisi tekrar oynatıldığında aynı yörünge
elde edilmez. Not: yalnızca başlangıç konumunun rastgele olması ortamı stokastik
yapmaz; stokastiklik geçiş ve gözlem gürültüsünden gelir.

\textbf{On-policy mi, off-policy mi?} PPO \textbf{on-policy}'dir. Veriyi üreten
davranış politikası ile iyileştirilen hedef politika aynıdır; her rollout, güncel
politika ile toplanır ve güncellemeden sonra atılır. Replay buffer yoktur. Bu, SAC ve
DQN'in (off-policy, replay buffer'lı) tersidir ve PPO'nun neden daha çok örnek
gerektirdiğini açıklar.

\textbf{Markov özelliği.} Durum, anlık lidar, yönelim, hız (önceki eylem yoluyla) ve
keşif istatistiklerini içerir. Tek bir gözlem, sonraki durumun dağılımını yaklaşık
olarak belirlemeye yeter; çünkü hareketli engellerin konumu adım sayısının
deterministik bir fonksiyonudur ve hız bilgisi gözlemde taşınır. Yine de gürültü ve
sınırlı görüş nedeniyle Markov özelliği yalnızca yaklaşık sağlanır.

\textbf{Kısmi gözlemlenebilirlik (POMDP).} Ortam \textbf{kısmi gözlemlenebilirdir}.
Gerçek durum ile gözlem ayrışır: (i) poz/yönelim gözleme odometri gürültüsüyle girer,
yani ajan gerçek konumunu tam bilmez; (ii) lidar yalnızca görüş hattındaki engelleri
görür, duvar arkası bilinmez; (iii) hareketli engellerin gelecekteki konumu doğrudan
gözlemde yoktur. Optimal karar için eksik olan bu bilgiyi ajan, önceki eylemi gözleme
ekleyerek ve politikanın hız/yön bağlamını taşımasıyla kısmen telafi eder.

\section{Sonuç}
PPO, replay buffer kullanmadan, on-policy bir politika gradyanı yöntemi olarak, aynı
ortamda hem rastgele hem sezgisel taban politikalarını belirgin biçimde geçmiştir.
Örnek verimliliği off-policy yöntemlere göre düşük olsa da, paralel ortam toplama ve
daha uzun eğitim ile yüksek oda kapsaması elde edilmiştir. Aynı ortam, seed ve
protokol SAC teslimiyle birebir korunduğundan, bu rapordaki sonuçlar ekipteki diğer
algoritmalarla doğrudan karşılaştırılabilir. Tüm grafikler \reff{sonuclar/loglar/}
altındaki ham loglardan yeniden üretilebilir.

\end{document}
